\chapter{Multi-Agent Systems and Agentic AI: Foundations}
\label{chap:background_mas}

\section{Introduction}

This chapter establishes the conceptual vocabulary used throughout the thesis. We first define agents and multi-agent systems, then review the coordination mechanisms most relevant to distributed surveillance---in particular task allocation by auction---and finally introduce the recent paradigm of \emph{agentic AI}, in which large language models act as reasoning engines inside agent architectures.

\section{Agents and Multi-Agent Systems}

\subsection{The Agent Abstraction}

Following Wooldridge \cite{wooldridge2009introduction}, an \textbf{agent} is a computer system situated in an environment, capable of \emph{autonomous action} in that environment in order to meet its design objectives. Four properties are commonly required:

\begin{itemize}
    \item \textbf{Autonomy}: the agent operates without direct external control of its internal state and decisions;
    \item \textbf{Reactivity}: it perceives its environment and responds in a timely fashion to changes;
    \item \textbf{Pro-activeness}: it exhibits goal-directed behavior, taking initiative rather than only reacting;
    \item \textbf{Social ability}: it interacts with other agents through some communication language or protocol.
\end{itemize}

In AURA-MAS, a camera perception agent exemplifies all four: it autonomously processes its own video stream (autonomy), emits events when zone rules fire (reactivity), continuously maintains tracking state and dwell-time counters (pro-activeness), and bids on verification tasks announced by a coordinator (social ability).

\subsection{Agent Architectures}

Classical agent architectures span a spectrum from purely \emph{reactive} (stimulus--response rules, subsumption) to \emph{deliberative} (symbolic planning over explicit world models), with the influential Belief--Desire--Intention (\gls{bdi}) model \cite{rao1995bdi} occupying a pragmatic middle ground: agents maintain \emph{beliefs} about the world, adopt \emph{desires} (goals), and commit to \emph{intentions} (plans). Hybrid layered architectures combine a fast reactive layer with a slower deliberative layer---a pattern that AURA-MAS mirrors at the system level: edge agents are predominantly reactive (frame-rate perception loops), while the coordination layer deliberates over aggregated hypotheses at a slower cadence, and the agentic explanation layer performs the slowest, most expensive reasoning only when justified.

\subsection{Multi-Agent Systems}

A \textbf{multi-agent system} (\gls{mas}) is a collection of interacting agents that jointly solve problems beyond the capability or knowledge of any individual agent \cite{wooldridge2009introduction}. \gls{mas} are attractive for surveillance for several structural reasons:

\begin{itemize}
    \item \textbf{Natural distribution}: sensors are physically distributed; assigning one agent per sensor matches the problem topology and keeps raw data local.
    \item \textbf{Scalability}: adding a camera means adding an agent, not re-dimensioning a central server.
    \item \textbf{Robustness}: the failure of one agent degrades coverage locally instead of collapsing the whole system.
    \item \textbf{Heterogeneity}: video, audio, and higher-level reasoning agents can coexist behind a common message interface.
\end{itemize}

Interaction is standardized by agent communication languages; the \gls{fipa} \gls{acl} defines performatives (\texttt{inform}, \texttt{request}, \texttt{cfp}, \texttt{propose}, \texttt{accept-proposal}, \ldots) whose semantics underpin most practical coordination protocols. AURA-MAS adopts the \emph{spirit} of \gls{fipa} interaction protocols while using lightweight \gls{json} messages over \gls{mqtt} topics, a common engineering compromise in modern \gls{iot}-oriented deployments \cite{light2017mosquitto}.

\section{Coordination and Task Allocation}

\subsection{The Coordination Problem}

Coordination is the process of managing dependencies between agents' activities. In distributed surveillance the canonical dependency is \emph{verification}: when one sensor produces an uncertain detection, which other sensor should spend resources to confirm or refute it? This is an instance of the \emph{multi-robot/multi-agent task allocation} problem, for which three broad families of solutions exist.

\subsection{The Contract Net Protocol and Auctions}

The Contract Net Protocol (\gls{cnp}) introduced by Smith \cite{smith1980contract} structures task allocation as a negotiation: a \emph{manager} announces a task (call for proposals), \emph{contractors} evaluate it against their local capabilities and submit \emph{bids}, and the manager awards the task to the best bidder. Market-based extensions formalize bids as utilities and prove strong properties for single-item auctions \cite{parker2002alliance}: allocation quality equals the best available bid, communication cost is linear in the number of bidders, and the mechanism degrades gracefully when agents fail to respond.

Auctions are particularly well suited to camera verification tasks because a camera's utility for verifying an event is inherently \emph{local, private knowledge}---its field-of-view overlap with the event zone, its current processing load, whether it originated the initial detection---which the bidder can compute cheaply and truthfully summarize as a scalar bid. AURA-MAS therefore adopts a single-round, sealed-bid auction (Section~\ref{sec:design_coordination}) with a fixed bidding window, and retains a round-robin scheduler as an ablation baseline.

\subsection{Alternative Approaches}

Two alternatives deserve mention. \emph{Centralized optimization} (e.g., Hungarian assignment over a global utility matrix) yields optimal allocations but requires the center to know all utilities, reintroducing the communication and single-point-of-failure costs that \gls{mas} avoid; recent work explores submodular optimization for multi-camera collaborative perception \cite{liu2026submodular}. \emph{Multi-agent reinforcement learning} (\gls{marl}) learns coordination policies end-to-end, with algorithms such as QMIX \cite{rashid2018qmix} and MADDPG \cite{lowe2017multiagent} demonstrating strong results in simulated cooperative sensing and patrolling \cite{zheng2026multiuav, yang2025multiagent}; standardized environments (PettingZoo \cite{terry2021pettingzoo}) and scalable training stacks (RLlib \cite{liang2018rllib}) exist. However, \gls{marl} policies are data-hungry, hard to certify, and opaque---undesirable properties for a safety- and legally-sensitive system. We therefore position \gls{marl}-based active sensing as future work (Chapter~\ref{chap:conclusion}) and keep the deployed coordination mechanism auditable.

\section{Agentic AI}

\subsection{From Language Models to Agents}

Large language models exhibit emergent capabilities for planning, tool use, and multi-step reasoning. The \emph{agentic AI} paradigm operationalizes these capabilities: an \gls{llm} is embedded in a control loop where it observes state, reasons (often via explicit ``thought'' traces as in ReAct \cite{yao2023react}), invokes tools, and iterates. Sapkota et al.\ \cite{sapkota2025ai} distinguish \emph{AI agents}---single-model, tool-using systems---from \emph{agentic AI}: orchestrated ensembles of specialized agents with shared memory and explicit coordination, which is precisely the regime of interest for surveillance, where perception specialists must feed reasoning specialists.

\subsection{Orchestration Frameworks}

A rich ecosystem of frameworks now supports multi-agent \gls{llm} orchestration: AutoGen structures multi-agent conversations \cite{wu2023autogen}; LangGraph models agent workflows as typed state graphs with explicit nodes, edges, and checkpointing \cite{langgraph2024}; CrewAI emphasizes role-based teams. For vision-centric pipelines, research systems such as VideoAgent \cite{fan2024videoagent} and Hawk \cite{cheng2025hawk} demonstrate \gls{llm}-orchestrated video understanding, while \gls{vlm}s such as GPT-4V \cite{openai2023gpt4}, Gemini \cite{google2023gemini}, Qwen-VL \cite{alibaba2023qwenvl}, and LLaVA \cite{liu2023visual} provide the underlying multimodal perception-to-language capability.

\subsection{Reliability and Guardrails}

The central obstacle to deploying generative models in safety-critical loops is \emph{hallucination}: fluent but unfounded output. Guardrail frameworks such as NeMo Guardrails \cite{rebedea2023nemo} interpose programmable checks between the model and the application. The design stance adopted in this thesis is stricter and architectural: the agentic layer is \emph{decision-decoupled}---it runs only after a deterministic policy engine has already made the alert decision---and \emph{evidence-grounded}---every evidence identifier cited in a generated report is mechanically checked against the event log, with automatic fallback to a template on violation. This turns the \gls{llm} from a potential safety liability into a strictly additive usability layer, a pattern we consider one of the transferable lessons of this work.

\section{Conclusion}

This chapter defined the agent abstraction, motivated the \gls{mas} architecture for distributed surveillance, selected auction-based task allocation as the coordination mechanism on grounds of efficiency and auditability, and framed agentic AI as a rule-guarded explanation layer rather than a decision maker. The next chapter reviews the perception techniques that edge agents rely upon.
